\section{Conclusion}\label{sec:conclusion}

Reporting is part of the visual task, not merely its formatting. When vision-language models report sets of objects, part of their error is in the report rather than in the seeing, and treating reporting as a chain from responses to records, sets and readouts makes that error analysable. The chain gives exact identities for what each operation does, a behavioural model of persistent errors, and a proof that outputs alone cannot certify correctness. It also separates three problems. Repair works where the information is already present: read the count from the model's list of records, not from the numbers it states, unless the list has degenerated; delete unconfirmed records only below the cost threshold; change the report structure, its decoding (declared aliases applied in code), the view or the model, rather than restating the same request more insistently. Detection needs a different source of evidence; on our click tasks a second model provided it better than a second call. Judging between reports needs labels or a configuration transferred from other data; three-call unanimity, chosen on five sources, improved set error on most over-reporting cells of the sixth and did not significantly worsen it in any of them. Every number in this paper can be regenerated from the released responses, with the DOTA, SKU-110K and FSC-147 labels rebuilt from their official downloads.
